\section[ПРОЕКТИРОВАНИЕ АРХИТЕКТУРЫ И МОДЕЛЕЙ СЕРВИСА]{Проектирование архитектуры и моделей сервиса}
\thispagestyle{plain}

В настоящей главе обоснован выбор архитектурного решения, описан
применяемый технологический стек и приведены формальные определения
моделей данных и алгоритмов сервиса EduQuiz: ролевой модели
пользователей, модели интерактивной учебной сессии, модели игровой
прогрессии, алгоритма индивидуального интервального повторения,
модели оценки качества тестового инструмента и модели данных.

\subsection{Обоснование выбора архитектурного решения}

Согласно требованиям, сформулированным в подразделе~1.3.3,
сервис должен размещаться в инфраструктуре, контролируемой
образовательной организацией или её региональным провайдером,
и предоставлять открытый программный интерфейс для интеграции
с информационными системами учебного заведения. Эти требования
исключают применение архитектурной модели Backend-as-a-Service
на зарубежной облачной платформе и определяют необходимость
перехода к классической трёхуровневой клиент-серверной архитектуре
с собственным серверным приложением.

В качестве альтернативных архитектурных моделей
рассматривались~\cite{fowler_poeaa,kleppmann2017}:
\begin{enumerate}
    \item монолитная архитектура с единым серверным приложением,
          обрабатывающим все запросы (преимущества~--- простота
          разработки и развёртывания, минимальные накладные расходы
          на сетевое взаимодействие; недостатки~--- ограниченная
          возможность независимого масштабирования отдельных
          подсистем);
    \item микросервисная архитектура с декомпозицией по бизнес-доменам
          (преимущества~--- независимое масштабирование, возможность
          применения различных технологий в разных сервисах;
          недостатки~--- высокая операционная сложность,
          избыточность для прогнозируемого профиля нагрузки
          учебной аудитории одного университета).
\end{enumerate}

С учётом профиля нагрузки разрабатываемого сервиса (десятки
параллельных учебных сессий по 30--500 участников каждая) выбрана
монолитная архитектура с собственным backend-приложением.
Описание архитектуры приведено в нотации модели
C4~\cite{hohpe2003}, обеспечивающей визуализацию системы
на четырёх уровнях детализации (системный контекст, контейнеры,
компоненты, программный код).

\subsection{Стек используемых технологий}

\textbf{Серверная часть.} В качестве языка программирования
серверной части выбран язык Go версии 1.25. Выбор обусловлен
сочетанием следующих характеристик: статической типизацией,
обеспечивающей обнаружение значительной части ошибок на этапе
компиляции; встроенной моделью конкурентности на основе горутин
и каналов, удобной для реализации сценариев с большим числом
одновременных сетевых соединений; низким расходом оперативной
памяти на одно сетевое соединение, что важно при поддержке
десятков WebSocket-соединений в рамках одной учебной сессии.
В качестве веб-фреймворка использован Gin~\cite{gin_docs};
для работы с базой данных~--- драйвер \texttt{pgx/v5}; для миграций
схемы~--- утилита \texttt{golang-migrate}; для поддержки WebSocket~---
библиотека \texttt{gorilla/websocket}~\cite{rfc6455}.

\textbf{Реляционная база данных.} В качестве реляционной системы
управления базами данных выбрана PostgreSQL версии
16~\cite{postgresql_docs}. Выбор обусловлен открытой лицензией,
поддержкой типа JSONB для хранения полуструктурированных данных
(варианты ответов на вопросы), наличием расширений \texttt{citext}
(регистронезависимые электронные адреса) и \texttt{pgcrypto}
(генерация UUID и хеширование), а также богатой экосистемой
средств администрирования. В качестве средства кеширования
и брокера событий между экземплярами серверного приложения
использована Redis версии 7. Обратный прокси Caddy
версии 2~\cite{caddy_docs} обеспечивает автоматический выпуск
TLS-сертификатов по протоколу ACME.

\textbf{Клиентская часть.} В качестве платформы разработки
клиентского приложения выбран Flutter версии 3.35~\cite{flutter_docs}
с языком программирования Dart. Выбор Flutter обоснован возможностью
разработки единой кодовой базы, разворачиваемой на четырёх целевых
платформах (Web, iOS, Android, macOS) без существенной потери
качества пользовательского интерфейса. Управление состоянием
приложения выполнено с использованием библиотеки
Riverpod~2.6~\cite{riverpod_docs}; маршрутизация~--- средствами
\texttt{go\_router~14} с декларативной защитой маршрутов
по роли пользователя.

\subsection{Архитектура клиент-серверного решения}

Архитектура разработанного сервиса представлена
на Рисунке~\ref{fig:c4-containers} в виде контейнерной диаграммы
C4. В состав системы входят следующие контейнеры:
\begin{enumerate}
    \item клиентское приложение (Flutter), функционирующее
          на устройствах преподавателя и обучающихся в четырёх
          вариантах сборки (Web, iOS, Android, macOS);
    \item серверное приложение (Go, Gin), реализующее бизнес-логику,
          аутентификацию и трансляцию событий интерактивных сессий;
    \item сервер базы данных (PostgreSQL~16), являющийся единственным
          источником истины для пользовательских данных, курсов,
          банков вопросов, состояния сессий и состояний игровой
          прогрессии;
    \item сервер кеширования (Redis~7), используемый для хранения
          refresh-токенов и как шина событий между экземплярами
          серверного приложения;
    \item обратный прокси-сервер (Caddy~2), являющийся точкой входа
          из публичной сети.
\end{enumerate}

\image{c4-containers.png}{Контейнерная диаграмма C4 сервиса EduQuiz}{fig:c4-containers}{0.95}

Доступ к серверу базы данных и серверу кеширования из публичной
сети закрыт; взаимодействие выполняется только в пределах внутренней
сети, что снижает поверхность атаки в соответствии с требованиями
OWASP ASVS~\cite{owasp_asvs}. Аутентификация пользователей
реализована на основе технологии JSON Web Token (JWT) в соответствии
со стандартом RFC~7519~\cite{rfc7519}; выпускается пара токенов~---
токен доступа со сроком действия 15 минут и токен обновления
со сроком действия 7 суток. Хеширование паролей выполнено
алгоритмом bcrypt с рабочим параметром $\text{cost} = 12$, что
обеспечивает время вычисления хеша порядка 250~мс на используемой
аппаратной конфигурации. Для предотвращения атак методом полного
перебора перед эндпоинтами аутентификации установлен ограничитель
скорости запросов по IP-адресу (token-bucket, не более 20 запросов
в минуту).

Передача событий интерактивных сессий между сервером и клиентами
выполнена по протоколу WebSocket в соответствии
с RFC~6455~\cite{rfc6455}. Размещение серверной инфраструктуры
на территории Российской Федерации соответствует требованиям
Федерального закона от 27.07.2006 № 152-ФЗ
«О персональных данных»~\cite{fz152}. Сервис развёрнут на домене
\texttt{eduquizz.ru} с автоматическим выпуском TLS-сертификата
по протоколу ACME.

\subsection{Ролевая модель пользователей}

В разработанной системе выделено три роли пользователей:
преподаватель, обучающийся, администратор. Распределение прав
по ролям определяется по принципу наименьших
привилегий~\cite{owasp_asvs,gost58833}: каждая роль обладает
минимальным набором прав, необходимым для выполнения целевых
сценариев.

Преподаватель создаёт курсы, формирует и редактирует банки вопросов,
проводит интерактивные учебные сессии и получает аналитические
отчёты по результатам сессий. Обучающийся подключается к сессии
по числовому коду, отвечает на вопросы, выполняет индивидуальное
повторение учебного материала и получает доступ к собственной
статистике. Администратор обладает полным доступом
к функциональности преподавателя и обучающегося, а также управляет
учётными записями и правами доступа к курсам.

Разграничение доступа к ресурсам реализовано на двух уровнях:
авторизация по принадлежности ресурса в обработчиках HTTP-запросов
и дублирующий контроль на уровне ограничений целостности базы
данных (внешние ключи и уникальные индексы). Двухуровневая защита
снижает риск нарушения доступа при ошибке в коде уровня приложения.
Сценарии использования системы представлены
на Рисунке~\ref{fig:usecase} в виде диаграммы прецедентов.

\image{usecase.png}{Диаграмма прецедентов сервиса EduQuiz}{fig:usecase}{0.95}

\subsection{Модель интерактивной учебной сессии}

Под интерактивной учебной сессией (далее~--- сессия) в разработанной
модели понимается ограниченный по времени сеанс коллективной работы
преподавателя и группы обучающихся, в ходе которого преподаватель
предъявляет последовательность вопросов из выбранного банка,
а участники предоставляют ответы через клиентское приложение.

Состояние сессии представлено конечным автоматом
$M = (S, S_0, F, T)$, где $S$~--- множество состояний;
$S_0 \in S$~--- начальное состояние; $F \subseteq S$~--- множество
финальных состояний; $T$~--- множество переходов. Множество
состояний:
$$S = \{\texttt{created}, \texttt{waiting}, \texttt{active}, \texttt{review}, \texttt{finished}\},$$
начальное состояние $S_0 = \texttt{created}$, множество финальных
состояний $F = \{\texttt{finished}\}$. Семантика состояний:
\begin{itemize}
    \item \texttt{created}~--- сессия создана преподавателем,
          числовой код выпущен, подключение участников ещё
          не открыто;
    \item \texttt{waiting}~--- открыто подключение участников,
          активный вопрос отсутствует;
    \item \texttt{active}~--- участникам предъявлен вопрос,
          ожидаются ответы;
    \item \texttt{review}~--- приём ответов на текущий вопрос
          завершён, участникам отображается правильный ответ
          и локальная статистика;
    \item \texttt{finished}~--- сессия завершена, результаты
          зафиксированы в базе данных.
\end{itemize}

Состояние сессии хранится на стороне сервера и является единым
источником истины; обновления транслируются клиентам по WebSocket.

В модели предусмотрены три режима проведения сессии, различающихся
способом синхронизации участников:
\begin{itemize}
    \item \texttt{classic}~--- синхронный групповой режим: все
          участники одновременно отвечают на общий текущий вопрос;
          темп задаёт преподаватель кнопкой «Дальше». Применим
          для семинарских занятий с обсуждением;
    \item \texttt{timer}~--- индивидуальный режим с ограничением
          времени: каждый участник проходит вопросы в собственном
          темпе с таймером на каждый вопрос; по истечении таймера
          ответ фиксируется как незасчитанный. Применим
          для контрольного тестирования;
    \item \texttt{race}~--- индивидуальный режим без ограничения
          времени с определением победителя~--- участника,
          первым правильно ответившего на все вопросы. Применим
          для соревновательного учебного формата.
\end{itemize}

Архитектурно режимы делятся на два класса: групповой режим
(\texttt{classic}) с общим состоянием на сессию и индивидуальные
режимы (\texttt{timer}, \texttt{race}) с состоянием, привязанным
к отдельному участнику.

\subsection{Модель игровой прогрессии}

Модель игровой прогрессии направлена на поддержание устойчивой
учебной мотивации обучающихся и повышение регулярности обращения
к учебному материалу за счёт механизмов внешней
мотивации~\cite{deterding2011,bartle1996}. Разработанная модель
включает четыре взаимосвязанных механизма: накопление очков опыта
(XP) за ответы в ходе сессии; систему уровней с замедляющимся
ростом; систему достижений (бейджей)~--- стимул к освоению
различных аспектов сервиса; учёт серий активных дней~--- стимул
к регулярному взаимодействию с сервисом между занятиями.

\textbf{Начисление очков опыта} за ответ на вопрос в ходе сессии
вычисляется по формуле~\eqref{eq:xp}.
\begin{equation}\label{eq:xp}
    \Delta XP = w_c \cdot d \cdot \left(1 - \alpha \cdot \frac{t}{T}\right) \cdot 10,
\end{equation}
где $w_c$~--- множитель за корректность ответа ($w_c = 1$
для правильного, $w_c = 0$ для неправильного); $d$~--- оценка
сложности вопроса (целое число от 1 до 5); $t$~--- время,
затраченное на ответ; $T$~--- максимальное время на ответ;
$\alpha = 0{,}5$~--- коэффициент штрафа за продолжительность
ответа. Множитель $10$ обеспечивает целочисленный диапазон значений
$\Delta XP$ в десятки единиц при сложности $d \in [1, 5]$,
что согласовано с пороговыми значениями уровней.

\textbf{Перевод накопленных очков опыта} в уровень $L$ выполняется
по формуле~\eqref{eq:level}.
\begin{equation}\label{eq:level}
    L = \left\lfloor \sqrt{XP / k_L} \right\rfloor, \quad k_L = 100.
\end{equation}
Квадратичная зависимость $XP(L) = k_L \cdot L^2$ обеспечивает
быстрый рост уровня в начале использования сервиса и замедление
по мере накопления опыта, что является классическим паттерном
ролевых игр.

\textbf{Система достижений} реализована как множество однократных
наград за выполнение определённых условий: получение первого
правильного ответа, прохождение сессии без ошибок, накопление
установленного объёма очков опыта, серия активных дней заданной
длины. Идемпотентность операции выдачи достижения обеспечивается
уникальным ключом \texttt{(user\_id, badge\_code, course\_id)}
в таблице выданных достижений.

\subsection{Алгоритм индивидуального интервального повторения}

Для индивидуального повторения учебного материала между сессиями
использован алгоритм SuperMemo SM-2~\cite{wozniak1990}, эффективность
которого в массовых системах распределённого обучения подтверждена
эмпирически.

Каждой паре «обучающийся, вопрос» сопоставляется состояние повторения
$\sigma = (EF, I, n)$, где $EF$~--- коэффициент лёгкости вопроса;
$I$~--- текущий интервал повторения в сутках; $n$~--- число
последовательных успешных повторений. Начальное состояние:
$EF_0 = 2{,}5$, $I_0 = 0$, $n_0 = 0$.

Обновление состояния после ответа обучающегося с оценкой качества
$q \in \{0, \ldots, 5\}$ выполняется по правилам~\eqref{eq:sm2-update}:
\begin{equation}\label{eq:sm2-update}
    \begin{cases}
        \text{если } q < 3: & n \leftarrow 0, \quad I \leftarrow 1; \\
        \text{если } q \geq 3 \text{ и } n = 0: & I \leftarrow 1, \quad n \leftarrow n + 1; \\
        \text{если } q \geq 3 \text{ и } n = 1: & I \leftarrow 6, \quad n \leftarrow n + 1; \\
        \text{если } q \geq 3 \text{ и } n > 1: & I \leftarrow \lceil I \cdot EF \rceil, \quad n \leftarrow n + 1.
    \end{cases}
\end{equation}

Коэффициент лёгкости $EF$ обновляется по формуле~\eqref{eq:sm2-ef}.
\begin{equation}\label{eq:sm2-ef}
    EF \leftarrow \max\bigl(1{,}3,\; EF + 0{,}1 - (5 - q) \cdot (0{,}08 + (5 - q) \cdot 0{,}02)\bigr).
\end{equation}
Дата следующего предъявления карточки определяется как
$d_{\text{next}} = d_{\text{cur}} + I$.

В классической реализации SM-2 оценка качества ответа $q$
запрашивается у обучающегося явно после предъявления правильного
варианта. В настоящей работе оценка $q$ формируется автоматически
по корректности ответа и времени реакции обучающегося
по правилу~\eqref{eq:q-score}.
\begin{equation}\label{eq:q-score}
    q =
    \begin{cases}
        5, & \text{ответ верный и } t \leq T/2; \\
        4, & \text{ответ верный и } T/2 < t \leq T; \\
        3, & \text{ответ верный и } t > T \text{ (или } T = 0\text{)}; \\
        2, & \text{ответ неверный и } t < T; \\
        1, & \text{ответ неверный, таймер истёк.}
    \end{cases}
\end{equation}

Автоматическое формирование $q$ устраняет необходимость отдельного
интерфейса самооценки и обеспечивает плавную интеграцию интервального
повторения в цикл интерактивной учебной сессии: после каждого
ответа в сессии помимо начисления очков опыта вызывается обновление
SM-2 состояния, что обеспечивает естественную аккумуляцию учебного
материала~--- один раз отвеченный вопрос автоматически попадает
в очередь повторения обучающегося. Адаптация алгоритма SM-2 для
совместной работы с интерактивной сессией составляет одну из задач
настоящей работы, поставленных в подразделе~1.4.

\subsection{Модель оценки качества тестового инструмента}

Помимо классических метрик успеваемости, для оценки качества
банка вопросов как измерительного инструмента применяются
психометрические показатели классической теории
тестов~\cite{cronbach1951,crocker_algina}.

\textbf{Коэффициент $\alpha$ Кронбаха} оценивает внутреннюю
согласованность тестового инструмента~--- тенденцию всех пунктов
измерять одну и ту же латентную способность; рассчитывается
по формуле~\eqref{eq:cronbach}.
\begin{equation}\label{eq:cronbach}
    \alpha = \frac{k}{k - 1} \left(1 - \frac{\sum_{i=1}^{k} \sigma_{Y_i}^2}{\sigma_X^2}\right),
\end{equation}
где $k$~--- число пунктов; $\sigma_{Y_i}^2$~--- дисперсия $i$-го пункта;
$\sigma_X^2$~--- дисперсия суммарного балла. Принято считать
тестовый инструмент надёжным при $\alpha \geq 0{,}7$.

\textbf{Качество отдельных пунктов} оценивается двумя показателями.
\emph{Индекс трудности} $p_i$ есть доля правильных ответов
на $i$-й вопрос по выборке участников. \emph{Точечно-бисериальная
корреляция} ответа на пункт с суммарным баллом по тесту
рассчитывается по формуле~\eqref{eq:rpb}.
\begin{equation}\label{eq:rpb}
    r_{pb,i} = \frac{\bar{X}_{+} - \bar{X}_{-}}{\sigma_X} \sqrt{p_i (1 - p_i)},
\end{equation}
где $\bar{X}_{+}$, $\bar{X}_{-}$~--- средние суммарные баллы
участников, соответственно правильно и неправильно ответивших
на $i$-й пункт. Положительные значения $r_{pb,i}$ свидетельствуют
о работающем пункте; отрицательные~--- о кандидате на исключение
или переформулировку.

Принятые в работе ориентиры интерпретации показателей:
$p_i \geq 0{,}85$~--- пункт слишком лёгкий; $p_i < 0{,}30$~---
слишком сложный; $r_{pb,i} \leq 0{,}2$~--- низкая дискриминативная
способность пункта. Расчёт показателей $\alpha$, $\{p_i\}$ и
$\{r_{pb,i}\}$ инкапсулирован в чистой функции серверной части
(\texttt{services.AnalyzeBank}), принимающей бинарную матрицу
ответов $U \in \{0, 1\}^{n \times k}$ и возвращающей сводный
отчёт. На клиенте отчёт отображается в форме психометрической
карточки с цветовой интерпретацией значения $\alpha$ и автоматическим
выделением пунктов, требующих доработки. Подробное описание
реализации психометрического анализатора приведено в подразделе~3.3
главы~3. Разработанная модель и её программная реализация составляют
второй из научных результатов настоящей работы.

\subsection{Модель данных}

База данных сервиса спроектирована как нормализованная (третья
нормальная форма) реляционная схема в PostgreSQL~16. Основные
сущности и их назначение приведены в табл.~\ref{tab:db-entities}.

\begin{table}[!htb]
\fontsize{11pt}{13pt}\selectfont
\centering
\caption{Основные сущности модели данных сервиса EduQuiz}
\label{tab:db-entities}
\smallskip
\begin{tabular}{|l|l|}
    \hline
    \textbf{Таблица} & \textbf{Назначение} \\
    \hline
        \texttt{users}            & учётные записи и роли пользователей \\
    \hline
        \texttt{courses}          & курсы, принадлежащие преподавателю \\
    \hline
        \texttt{question\_banks}  & банки вопросов в составе курса \\
    \hline
        \texttt{questions}        & вопросы с вариантами ответа в формате JSONB \\
    \hline
        \texttt{rooms}            & интерактивные учебные сессии \\
    \hline
        \texttt{participants}     & участники сессий \\
    \hline
        \texttt{answers}          & ответы участников на вопросы сессии \\
    \hline
        \texttt{xp\_log}          & журнал начислений очков опыта \\
    \hline
        \texttt{user\_progress}   & агрегат игровой прогрессии (XP, уровень, серии) \\
    \hline
        \texttt{badges}, \texttt{user\_badges} & справочник и выданные достижения \\
    \hline
        \texttt{sm2\_states}      & состояния интервального повторения по SM-2 \\
    \hline
\end{tabular}
\end{table}

Полные определения таблиц и индексов содержатся в файлах миграций
\texttt{migrations/0001\_init}~--- \texttt{migrations/0010}.
Ключевые конструктивные решения модели данных описаны ниже.

Уникальное ограничение
$\texttt{UNIQUE} (\texttt{room\_id}, \texttt{question\_id}, \texttt{participant\_id})$
на таблице \texttt{answers} предотвращает повторную регистрацию
ответа одного участника на один и тот же вопрос в рамках одной
сессии и обеспечивает устойчивость системы к повторным отправкам,
вызванным двойным нажатием или сетевыми сбоями.

Тип вопроса описывается ENUM
$\texttt{question\_kind} = \{\texttt{single\_choice}, \texttt{multi\_choice}, \texttt{open\_text}, \texttt{rating}, \texttt{qna}, \texttt{true\_false}, \texttt{numerical}\}$.
Варианты ответа и правильный ответ хранятся в JSONB-полях,
что обеспечивает гибкость без необходимости менять схему таблицы
при появлении новых типов.

Конечный автомат состояния сессии (подраздел~2.5) реализован через
ENUM \texttt{room\_status}; список заданных вопросов и очередь
прохождения хранятся в массивах \texttt{question\_order UUID[]}
и \texttt{asked\_question\_ids UUID[]} таблицы \texttt{rooms}.

Состояние интервального повторения для пары
$(\texttt{user\_id}, \texttt{question\_id})$ хранится в таблице
\texttt{sm2\_states} с полями \texttt{ef REAL}, \texttt{interval\_days
INTEGER}, \texttt{repetitions INTEGER}, \texttt{due\_date DATE}.
Индекс $\texttt{idx\_sm2\_due} (\texttt{user\_id}, \texttt{due\_date})$
обеспечивает быструю выборку карточек, запланированных
на текущую дату, при росте банка SM-2-состояний до десятков
тысяч записей.

Денормализованный агрегат \texttt{user\_progress} с первичным ключом
$(\texttt{user\_id}, \texttt{course\_id})$ и полями
\texttt{total\_xp INTEGER}, \texttt{level INTEGER},
\texttt{streak\_days INTEGER} мотивирован требованием обновлять
таблицу лидеров за время не более 100~мс при $N = 100$ одновременно
отвечающих участниках. Согласованность с журналом \texttt{xp\_log}
обеспечивается одновременным обновлением обеих таблиц в одной
транзакции.

\textbf{Выводы по главе 2.}
В главе обоснован выбор монолитной клиент-серверной архитектуры
сервиса EduQuiz и выполнено её описание в нотации C4. Обоснован
выбор стека программных средств: Go~1.25 с фреймворком Gin
для серверной части, PostgreSQL~16 и Redis~7 для хранения данных
и кеширования, Caddy~2 в качестве обратного прокси, Flutter~3.35
для кроссплатформенного клиентского приложения. Разработана
ролевая модель пользователей с разграничением доступа на двух
уровнях (приложение и база данных). Разработана модель интерактивной
учебной сессии в виде конечного автомата с поддержкой трёх режимов
проведения. Разработана модель игровой прогрессии с формулами
расчёта очков опыта и уровней. Адаптирован алгоритм SuperMemo SM-2
для режима индивидуального повторения, объединённого с интерактивной
сессией через автоматическое формирование оценки качества ответа
из пары «корректность~--- время». Разработана модель оценки качества
тестового инструмента средствами классической теории тестов
($\alpha$ Кронбаха, индекс трудности, точечно-бисериальная
корреляция). Спроектирована модель данных в форме нормализованной
реляционной схемы PostgreSQL.

\newpage
